\section{Genetic Algorithm Hyperparameter Optimization}

The performance of Deep Reinforcement Learning algorithms is strongly
influenced by the choice of hyperparameters. Parameters such as learning
rate, rollout length, batch size, discount factor, replay buffer size, and
entropy regularization directly affect training stability, convergence speed,
and final trading performance. To automate the hyperparameter selection
process, a \textbf{Genetic Algorithm (GA)} was employed as an evolutionary
optimization method.

In the proposed formulation, each candidate hyperparameter configuration is
represented as a \textbf{chromosome}. A chromosome consists of multiple
\textbf{genes}, where each gene corresponds to a specific hyperparameter.
For example, a PPO chromosome contains genes representing the learning rate,
rollout length, batch size, discount factor, generalized advantage estimation
parameter, entropy coefficient, and clipping range. Similarly, chromosomes for
off-policy algorithms such as DDPG, TD3, and SAC contain genes describing
replay buffer sizes, target network update parameters, and other
algorithm-specific hyperparameters.

The optimization process begins by generating an initial
\textbf{population} of chromosomes through random sampling of the predefined
search space. Each chromosome is then evaluated by training a DRL agent using
the corresponding hyperparameter configuration for a fixed number of
timesteps. After training, the resulting agent is evaluated in the trading
environment, and its performance is used to compute a \textbf{fitness}
score. In this work, the fitness function is based on the final portfolio
value achieved during evaluation, providing a direct measure of trading
performance.

Following fitness evaluation, the population is ranked according to fitness
score and a \textbf{selection} process is applied. Selection determines which
chromosomes are allowed to contribute genetic material to the next generation,
favoring individuals with higher fitness while discarding poorly performing
solutions. In the proposed implementation, selection is performed through an
elite-based strategy in which the highest-ranking chromosomes form a candidate
pool for reproduction.

The strongest individuals are additionally preserved through
\textbf{elitism}, ensuring that the best solutions discovered so far are
carried unchanged into the next generation. New chromosomes are subsequently
generated through \textbf{crossover} and \textbf{mutation} operations.

The crossover operator combines genes from two parent chromosomes to produce
offspring that inherit characteristics from both solutions. More specifically,
a uniform crossover strategy is employed, where each gene is independently
inherited from either parent with equal probability. This allows successful
hyperparameter combinations to propagate throughout the population. Mutation
introduces random modifications to individual genes with a predefined
probability, increasing genetic diversity and enabling the search process to
explore previously unseen regions of the hyperparameter space.

The cycle of evaluation, selection, crossover, and mutation is repeated over
multiple \textbf{generations}. As evolution progresses, poorly performing
configurations are gradually eliminated, while high-fitness chromosomes become
increasingly dominant within the population. The best chromosome encountered
during the evolutionary process is selected as the final hyperparameter
configuration for the corresponding DRL algorithm.

Compared to manual tuning or exhaustive grid search, the Genetic Algorithm
provides a population-based search strategy capable of efficiently exploring
high-dimensional and highly non-linear hyperparameter spaces. By balancing
exploration through mutation and exploitation through selection and elitism,
the GA is able to discover competitive hyperparameter configurations while
requiring substantially fewer evaluations than exhaustive search methods,
making it particularly suitable for computationally expensive financial
reinforcement learning environments.